\subsection{Extension of scalars in an integrally closed algebra}

PROPOSITION 19. Let \(k\) be a field, \(L\) a separable extension of \(k\) and \(R\) an integrally closed \(k\) -algebra. If the ring \(L \otimes_k R\) is an integral domain, it is integrally closed.

Let K be the field of fractions of R; as k is a field, \(L \otimes_{k} R\) is canonically identified with a sub-k-algebra of \(L \otimes_{k} K\) and L and R with sub-k-algebras of \(L \otimes_{k} R\) . Moreover, since an element \(s \neq 0\) of R is not a divisor of 0 in R, \(1 \otimes s\) is not a divisor of zero in \(L \otimes_{k} R\) since L is flat over k (Chapter I, §2, no. 3); identifying s with \(1 \otimes s\) , it is therefore seen that, if S = R - \{0\}, \(L \otimes_{k} K\) is identified with \(S^{-1}(L \otimes_{k} K)\) ; as \(L \otimes_{k} R\) is assumed to be an integral domain, \(L \otimes_{k} K\) is thus identified with a subring of the field of fractions \(\Omega\) of \(L \otimes_{k} R\) .

\begin{enumerate}
\def\labelenumi{(\arabic{enumi})}
\tightlist
\item
  Suppose that \(\mathbf{L}\) is a finite extension of \(k\) ; then \(\mathbf{L} \otimes_k \mathbf{K}\) is an algebra of finite rank over \(\mathbf{K}\) and by hypothesis has no divisor of 0; hence it is a field (Algebra, Chapter V, §2, no. 1, Proposition 1) and therefore it is in this case the field of fractions \(\Omega\) of \(\mathbf{L} \otimes_k \mathbf{R}\) . Let \((w_1, \ldots, w_n)\) be a basis of \(\mathbf{L}\) over \(k\) , which is therefore also a basis of \(\mathbf{L} \otimes_k \mathbf{K}\) over \(\mathbf{K}\) . There exists a basis \((w_1^*, \ldots, w_n^*)\) of \(\mathbf{L}\) such that \(\mathrm{Tr}_{\mathrm{L/k}}(w_i w_j^*) = \delta_{ij}\) (no. 6, Lemma 3); every \(z \in \mathbf{L} \otimes_k \mathbf{K}\) may be written uniquely \(z = \sum_{i=1}^{n} a_i w_i\) where \(a_i \in \mathbf{K}\) ; then
\end{enumerate}

\[
\operatorname{Tr} _ {(\mathrm{L} \otimes \mathrm{K}) / \mathrm{K}} (z w _ {j} ^ {*}) = \sum_ {i = 1} ^ {n} a _ {i} \operatorname{Tr} _ {(\mathrm{L} \otimes \mathrm{K}) / \mathrm{K}} (w _ {i} w _ {j} ^ {*})
\]

and as in L the traces \(\mathrm{Tr}_{(L\otimes K)/K}\) and \(\mathrm{Tr}_{L/K}\) coincide (Algebra, Chapter VIII, § 12, no. 2, formula (13)) finally \(\mathrm{Tr}_{(L\otimes K)/K}(zw_j^*) = a_j\) for \(1 \leqslant j \leqslant n\) . Note on the other hand that the elements of L are integral over \(k\) and hence also over R (no. 1, Corollary 1 to Proposition 2); therefore (no. 1, Proposition 5) \(\mathbf{L} \otimes_k \mathbf{R}\) is integral over R. This being so, suppose that \(z \in \mathbf{L} \otimes_k \mathbf{K}\) is integral over \(\mathbf{L} \otimes_k \mathbf{R}\) ; then \(z\) is also integral over R (no. 1, Proposition 6), hence so is \(zw_j^*\) and therefore also \(a_j = \mathrm{Tr}_{(L\otimes K)/K}(zw_j^*)\) for \(1 \leqslant j \leqslant n\) (no. 6, Corollary 2 to Proposition 17). As R is integrally closed, \(a_j \in \mathbf{R}\) for all \(j\) and hence \(z \in \mathbf{L} \otimes_k \mathbf{R}\) , which proves the proposition in this case.

\begin{enumerate}
\def\labelenumi{(\arabic{enumi})}
\setcounter{enumi}{1}
\tightlist
\item
  Suppose now that \(\mathbf{L}\) is a finitely generated separable extension of \(k\) ; then there exists a separating transcendence basis \((x_1, \ldots, x_d)\) of \(\mathbf{L}\) over \(k\) (Algebra, Chapter V, § 9, no. 3, Theorem 2); as \(\mathbf{L}\) and \(\mathbf{K}\) are algebraically disjoint over \(k\) in the field \(\Omega\) (Algebra, Chapter V, § 5, no. 4), the \(x_i\) are algebraically independent over \(\mathbf{K}\) ; hence \(\mathrm{R}[x_1, \ldots, x_d]\) is integrally closed (no. 3, Corollary 2 to Proposition 13). Let \(\mathbf{T}\) be the set of elements \(\neq 0\) of the ring \(\mathbf{A} = k[x_1, \ldots, x_d] \subset \mathbf{L}\) , so that the field \(k_1 = k(x_1, \ldots, x_d) \subset \mathbf{L}\) is equal to \(\mathbf{T}^{-1} k[x_1, \ldots, x_d]\) ; then
\end{enumerate}

\[
\begin{array}{r l} k _ {1} \otimes_ {k} \mathrm{R} & = (\mathrm{T} ^ {- 1} \mathrm{A}) \otimes_ {k} \mathrm{R} = \mathrm{T} ^ {- 1} \mathrm{A} \otimes_ {\mathrm{A}} (\mathrm{A} \otimes_ {k} \mathrm{R}) = \mathrm{T} ^ {- 1} (\mathrm{A} \otimes_ {k} \mathrm{R}) \\ & = \mathrm{T} ^ {- 1} \mathrm{R} [ x _ {1}, \dots , x _ {d} ] \end{array}
\]

by the associativity of the tensor product, hence this domain is integrally closed (no. 5, Corollary 1 to Proposition 16). But \(\mathbf{L} \otimes_k \mathbf{R}\) is identified with \(\mathbf{L} \otimes_{k_1} (k_1 \otimes_k \mathbf{R})\) and by definition \(\mathbf{L}\) is a finite separable extension of \(k_1\) ; it follows therefore from (1) that \(\mathbf{L} \otimes_k \mathbf{R}\) is integrally closed.

\begin{enumerate}
\def\labelenumi{(\arabic{enumi})}
\setcounter{enumi}{2}
\tightlist
\item
  General case. If z is an element of \(\Omega\) which is integral over \(L \otimes_{k} R\) , it satisfies a relation of the form \(z^{m} + b_{1}z^{m-1} + \cdots + b_{m} = 0\) , where the \(b_{i}\) belong to \(L \otimes_{k} R\) ; then there exists a finitely generated sub-extension \(L'\) of L over k such that the \(b_{i}\) belong to \(L' \otimes_{k} R\) for \(1 \leqslant i \leqslant m\) and z to \(L' \otimes_{k} K\) . Then it follows from (2) that \(z \in L' \otimes_{k} R\) and hence \(L \otimes_{k} R\) is integrally closed.
\end{enumerate}

* Let V be an affine irreducible algebraic variety, k a defining field of V and R the ring of functions regular on V defined over k; if R is integrally closed, V is called normal over k; Proposition 19 shows that, if V is normal over k, it remains normal over every separable extension L of k. *

COROLLARY. Let \(k\) be a field and \(\mathbf{R}\) and \(\mathbf{S}\) two integrally closed \(k\) -algebras. Suppose that the ring \(\mathbf{R} \otimes_k \mathbf{S}\) is an integral domain and that the fields of fractions \(\mathbf{K}\) and \(\mathbf{L}\) of \(\mathbf{R}\) and \(\mathbf{S}\) respectively are separable over \(k\) . Then the ring \(\mathbf{R} \otimes_k \mathbf{S}\) is an integrally closed domain.

As R and S are identified with subalgebras of \(R \otimes_{k} S\) , K and L are identified with subfields of the field of fractions \(\Omega\) of \(R \otimes_{k} S\) which are linearly disjoint over k (Algebra, Chapter V, §2, no. 3, Proposition 5). It then follows from Proposition 19 that \(R \otimes_{k} L\) and \(K \otimes_{k} S\) are integrally closed domains; as their intersection is \(R \otimes_{k} S\) (Chapter I, §2, no. 6, Proposition 7), \(R \otimes_{k} S\) is an integrally closed domain (no. 2, Corollary to Proposition 8).

* Given two irreducible affine varieties V, W defined over k, their product \(V \times W\) is an affine variety and the ring of functions regular on \(V \times W\) is identified with the tensor product over k of the ring of functions regular on V and the ring of functions regular on W. The Corollary to Proposition 19 shows that, if V and W are normal over k, then \(V \times W\) is normal over k. *

